\chapter{Design Analysis}

The design analysis of the electromagnetic brake joint is divided into four hierarchical layers, including the WSoI which encompasses the testing enviornment and the PCMA arm, the NSoI representing the joint, its functional subsystems, and the individual components. The physical development of the joint was iterative, however, the resulting architecture is documented in a logical, top-down fashion instead of a chronological sequence. To describe the system's integration, the analysis is split into three domains which constitute the mechanical design, the electronic design, and the software design. The mechanical design outlines the physical configuration and loading paths. The electronic design describes the peripheral circuitry and power interfacing. The software design defines the control logic and the state machine controlling the electromagnetic brake.

\section{Mechanical Design}

Figure~\ref{fig:screenshot-2026-10-21-103108} below, shows the 3D CAD assembly of the electromagnetic brake produced as part of the design analysis after integration and refinement of the individual subsystem concepts.
\begin{figure}[H]
	\centering
	\includegraphics[width=0.7\linewidth]{"../../../../../../../../../Downloads/Screenshot 2026-10-21 103108"}
	\caption{Electromagnetic Brake Assembly}
	\label{fig:screenshot-2026-10-21-103108}
\end{figure}
The top-level system elements are listed in the drawing's materials list and are marked with balloon numbers in Figure~\ref{fig:screenshot-2026-10-21-103108}. The different views and dimensions of the engineering drawing depict the position of physical interfaces with elements external to the NSoI and the spatial constraints for the respective system elements inside the NSoI.
The overall configuration uses an electromagnetic brake mounted coaxially on the joint shaft. In the configuration, the base link and the output link are integrated with the internal actuation and sensing components. The joint is supported by two RS PRO \SI{20}{mm} ID 2-Row ball bearings mounted with an inteference fit and secured by RS PRO \SI{32}{mm} OD one-piece steel shaft collars to prevent axial slipping. The bearing housing is integrated directly into the stationary joint housing. The housing encases the internal mechanisms and is held rigid by the base link. On one side of the shaft, the GTS06-OC-RA1000B-2M optical encoder is fitted into a centered hole in the shaft and secured with a scrub screw to the flat surface of the encoder's shaft\footnote{See Appendix \#\# for more detailed drawings regarding the joint shaft  design.}. On the opposite end, the electromagnetic brake assembly interfaces directly with the shaft. The electromagnetic brake consist of solenoids, friction plate (laser-cut by Fabrinox), and a brake disc. All other components are fabricated from piping and flat bars that are welded and/or fastened together.

The load path originates at the weight of the articulation link, which is transfered as torque through the shaft, and is ultimately countered by the electromagnetic force required to hold the link in position, when the electromaget's power is shut off.

Although TPM\_03 was already determined, the weight of the articulating link, accounting for the probe hub and probe, were calculated under Equation~\ref{eq:Articulation_link_weight} below, to verify if the TPM would still hold. Notes, a density of \SI{2700}{kg/m^3} was assumed for aluminium and the same safety factor and FBD under Section 3.4.3 was used\citep{Hibbler:2017}.
\begin{equation}
	\label{eq:Articulation_link_weight}
	\begin{aligned}
		m_{\text{link}} &= \rho_{\text{Al}} \cdot V_{\text{total}} \\
		&= \rho_{\text{Al}} \cdot (V_{\text{pipe}} + V_{\text{probe hub}} + V_{\text{probe}}) \\
		&= \rho_{\text{Al}} \cdot \left[ \underbrace{\pi (r_o^2 - r_i^2) \cdot L_{\text{pipe}}}_{\text{pipe volume}} + \underbrace{\pi(r_{\text{top}}^2 + r_{\text{bottom}}^2) \cdot L_{\text{probe hub}}}_{\text{probe hub volume}} + \underbrace{\pi r_{\text{probe}}^2 \cdot L_{\text{probe}}}_{\text{probe volume}} \right] \\
		&= 2700 \cdot \Big[ \pi(0.016^2 - 0.013^2)(0.55) \\
		&\quad + \pi(0.017^2 + 0.013^2)(0.01) + \pi(0.005^2)(0.1) \Big] \\
		&= \SI{0.4659}{kg} \\[1ex]
		\therefore T_{\text{holding}} &= \text{SF} \cdot m_{\text{link}} \cdot g \cdot l_{\text{link}} \\
		&= (1.5)(0.4659)(9.81)(0.325) \\
		&= \SI{2.228}{N.m}
	\end{aligned}
\end{equation}
This Equation~\ref{eq:Articulation_link_weight} confirms that the TPM is valid, however, the load path is still determined using the \SI{7.725}{Nm}. Next step was to size the dimensions of the brake in order to accurately determine how much force it will need to produce. An arbitrary diameter of \SI{0.1}{m} was chosen was chosen for the brake disc and friction plate, as discussed in Section 2.4.1. This means that the force required to be actuated by the brake can be calculated using Equation~\ref{eq:F_holding} below:
\begin{equation}
	\label{eq:F_holding}
	F_{\text{}}
\end{equation}
The friction plate was chosen to be mild steel and the brake disc to be wood. Their mating surfaces will be assumed to produce a friction coefficient of $\mu=0.6$

This configuration was chosen because it provides a reliable, uniform clamping force that minimises parasitic deflections and preserves the run-out tolerances necessary for the PCMA's spatial accuracy. Addionally, by housing the brake internally, the design is able to achieve the full \SI{90}{\degree} operational articulation sweep.

To maintain structural integrity and positional stability during locking, the electromagnetic brake was sized to deliver a static holding of \SI{7.725}{Nm}, corrsponding to a electromagntive force of \SI{3 550}{N}. The freebody diagrams (FBDs) and design calculations supporting the mechanical operating conditions of the joint under Appendix \#\#.

\section{Electronic Design}

The purpose of the switching circuit is to bridge the significant power gap between the high-power electromagnet and the low-power control logic, as well the data management system. The stm32 Nucleo microcontroller outputs a low-current control signal, which is insufficient to directly drive the electromagnetic brake that requires a \SI{12}{A} operating current at \SI{24}{V_DC}. This means that peripheral circuitry is required to safely step down power for the logic and allow the microcontroller to switch the electromagnet without sustaining damage. Figure~\ref{fig:wiring-diagram} below shows a wiring diagram detailing these connections.

\begin{figure}[h]
	\centering
	\includegraphics[width=0.7\linewidth]{"../../../../../../../../../Downloads/Wiring Diagram"}
	\caption{Wiring Diagram}
	\label{fig:wiring-diagram}
\end{figure}

To power the logic components, an XL7015 step-down buck converter is implemented. Its purpose is to reduce the \SI{24}{V} supply to a safe operating voltage suitable for the STM32 microcontroller and the rotary encoder.

To control the energisation of the brake, an Infineon IRF640NPBF N-channel metal-oxide-semiconductor field-effect transistor (MOSFET) is utilised as a low-side switch. The MOSFET acts as a digital relay. It's purpose is to take the control signal from the microcontroller\footnote{From the microcontroller, one of the pins is configured as a general pin input/output, which outputs control signal typically between \SI{20}{mA} and \SI{25}{mA} at an operating voltage level of \SI{3.3}{V}.} and use it to safely turn the high-current electromagnet on and off.

There are supporting components in the switching circuit which are critical and protective of the MOSFET, including the \SI{330}{\Omega} gate resistor, the \SI{10}{k\Omega} pull-down resistor, and the TO-220 flyback diode. The gate resistor is positioned between the microcontroller's signal pin and MOSFET gate. It's used to limit the sudden current flow needed to charge the MOSFET's internal capacitance in the event of switching the electromagnet on or off. This will prevent the microcontroller pin from burning out. The pull-down resistor is connected between MOSFET gate and drain. This allows the pull-down resistor to force the gate voltage to zero whenever the microcontroller is powering up, resetting, or outputting a floating signal. As a result, the MOSFET defaults to an "off" state, keeping the joint safely locked and preventing unintentional energisation. The flyback diode is placed in parallel across the electromagnetic coil. Its purpose is to minimise voltage spikes that could damage the MOSFET. When the electromagnet's current is suddenly cut off, the collapsing magnetic field generates a back electromotive force (EMF) due to its inductance. The back-emf would result is a high-voltage spike that exceeds the MOSFET's breakdown limit cause catastrophic failure. The diode provides a safe, low resistance path to dissipate the voltage spike.

For the sensing subsystem, the GTS06-OC-RA1000B optical rotary encoder is connected into the logic power rails, with its signal lines that were connected directly to the microcontroller to transmit real-time angular displacement data. The associated electrical calculation sizing these components and evaluating the MOSFET's thermal dissipation are detailed in Appendix \#\#.

\section{Software Design}

In mechatronics systems, such as the electromagnetic brake system, it is necessary to define the interactions between actuators (electrmoangetic brake), sensors (encoder), and the control software for reliable operation. The software architecture for the brake system uses a supervisory control and data acquisition (SCADA) hierachy to manage real-time monitoring and actuation processes. This hierarchy consists of the following levels, as shown in Figure~\ref{fig:softwareflowchart}:

\begin{itemize}
	\item Level 0 (field devices) consists of the hardware interacting directly with the optical rotary encoder, the IRF640N MOSFET, and the electromagnet brake.
	\item Level 1 (control logic) is executed by the STM32 Nucleo microcontroller, which acts as input/output processing module.
	\item Level 2 (supervisory) is handled by the connected computer that recieves angular displacement data from the microcontroller for logging and visualisation.
\end{itemize}

To fulfill the system's FRs, the microcontroller's pinout is configured to interface with the Level 0 devices. The rotary encoder's Phase A and Phase B signal lines are connected to pins PA0 and PA1. These pins are routed internally to the STM32's hardware timer (TIM2). The encoder outputs a quadrature pulse\footnote{A quadrature pulse is a signal that consists of two identical square waves offset by \SI{90}{\degree}.}, which allows the system to determine rotational position and direction of the shaft based on which phase leads the other. By utilising X4 encoding, which is when pulses are counting on the rising and falling edge, for both channels, the 1000 pulse per revolution (PPR) encoder yields 4000 distinct counts per 360-degree rotation.

By operating TIM2 in Encoder Mode, the microcontroller is able to use its dedicated hardware circuitry to autonomously track the high-frequency quadrature pulses and update a 32-bit memory register entirely in the background. This type of tracking eliminates the need for CPU-intensive interrupts, which would force the processor to halt execution of the code to count thousands of pulses per second. As a result, the CPU is freed up to properly manage other tasks. Additionally, the MOSFET gate is driven by pin PC2, which is configured as a digital output. Pin PC13 is configured as a hardware interrupt (EXTI), which is tied to the user input button. Lastly, pins PA2 and PA3 are configured for USART2 serial transmission to communicate the level 2  supervisory computer.

The control strategy implemented for the electromagnetic brake is an open-loop, bang-bang controller. Due to the binary mechanical requirement\footnote{The joint must be locked in placed, or released, free to rotate.}, duty cycle control is unnecessary. As shown in Figure~\ref{fig:softwareflowchart}, the system initialises into a locked state by defaulting PC2 low, thereby creating a fail-safe protocol. 

\begin{figure}[H]
	\centering
	\includegraphics[width=0.52\linewidth]{../../../../../../../../../Downloads/Software_Flowchart}
	\caption{Software Design of the PCMA Operation}
	\label{fig:softwareflowchart}
\end{figure}

The main loop reads the 32-bit counter from TIM2, converts the raw pulse count into an angular displacement into degrees, and transmits this data at \SI{20}{Hz} to the supervisor computer via USART2.

To manage the energisation and de-energisation of the electromagnet, the software uses the external interrupt on PC13. When the user presses the button, the interrupt service routine (ISR) pauses the main looop and toggles the joint's variable (joint\_locked) state. If the variable goes high (joint\_locked == 1), the controller drives the MOSFET gate high (\SI{3.3}{V}). This energises the solenoid and opens the air gap. Due to the asynchronous execution of the interrupt, the user does not experience lag when manually operating the arm. This is done to ensure a seamless human machine interface (HMI) between the mechanical operation and the software control state.
